
\section{H1预告更新与Q1利润桥}
公司7月14日盘后披露2026H1归母净利润29-30亿元、扣非24-25亿元、EPS1.58-1.64元。扣除Q1实际归母11.10亿元后，Q2归母净利润约17.90-18.90亿元，较Q1环比约61%-70%。这验证了公司层面的盈利提速，但不等于光模块分部利润已经单独披露。\n\n2026Q1收入同比增长52.72\%，归母净利润同比增长143.47\%，毛利率同比提升5.19个百分点。公司解释包括索尔思和GMD新增并表、光模块客户订单加急及传统业务稳定；H1预告进一步指出光模块整合效应、新增产能爬坡、新客户导入和数据中心相关投资收益开始贡献。通过公司披露的索尔思占比反推，Q1索尔思收入约为21.05亿元、利润约为5.95亿元；这两个绝对数是“公司披露比例乘合并报表”的推算，不是公司单独公布的分部利润表。

\begin{exhibitbox}[表4-1\quad 2026E盈利桥（House model）]
\begin{tabularx}{\exhibitboxwidth}{L{3.1cm}R{1.7cm}R{1.7cm}R{1.8cm}X}
\textbf{业务} & \textbf{收入} & \textbf{净利率} & \textbf{净利} & \textbf{证据与边界} \\
传统FPC/PCB & 260 & 4.50\% & 11.70 & 2025收入和Q1稳定出货，正常制造倍数 \\
AI PCB增量 & 60 & 10.00\% & 6.00 & 能力与扩产已知，收入/ASP/利用率为假设 \\
光模块/光芯片 & 280 & 17.00\% & 47.60 & Q1贡献与2025毛利已知，未来为条件信用 \\
显示模组 & 60 & 1.50\% & 0.90 & 低毛利稳定业务 \\
精密组件/GMD & 80 & 2.50\% & 2.00 & 平台扩张，利润率保守 \\
其他 & 5 & 3.00\% & 0.15 & 残余项 \\
\textbf{合计} & \textbf{745} & — & \textbf{68.35} & 四舍五入后为2026E归母净利CNY6.80bn \\
\end{tabularx}
\sourcenote{analysis/growth\_earnings\_model.md；data/growth\_driver\_model.json；金额单位均为CNY100m。}
\end{exhibitbox}

\section{现金流与资本强度}
2025年经营现金流53.07亿元、投资现金流净流出82.83亿元；2026Q1经营现金流11.27亿元、投资现金流净流出25.28亿元、筹资现金流净流入24.01亿元。Q1在建工程同比上升46.96\%，并购贷款和融资成本增加。增长模型必须同时承认利润弹性和现金流压力，不能只引用利润增速。

\section{当前价格隐含增长}
以260.37元和18.316亿股计算，公司市值约4768.96亿元。H1预告中点29.5亿元完成House 2026E归母净利润68亿元的43.4%，剩余H2仍需38.5亿元；若以30倍2027E PE解释现价，需要每股收益8.68元、归母净利润约159.0亿元；本报告House 2027E归母净利润132.1亿元，差额约26.9亿元。这个差额就是价格对更快光模块转化、更高AI PCB利润和更长估值久期的要求。

\begin{exhibitbox}[表4-2\quad 需要持续验证的驱动]
\begin{tabularx}{\exhibitboxwidth}{L{3.0cm}X L{3.0cm}X}
\textbf{驱动} & \textbf{基准假设} & \textbf{确认方式} & \textbf{失效条件} \\
光模块收入 & 2026E CNY28.0bn、2027E CNY50.0bn & 正式半年报分部收入和毛利 & Q2光模块收入低于CNY2.4bn \\
光模块利润率 & 2026E净利率17\% & 分部毛利、费用和现金流 & 连续两季低于15\% \\
AI PCB & 2026E增量收入CNY6.0bn & 高端产品收入、ASP、良率和利用率 & 仅有扩产无收入里程碑 \\
资金效率 & 经营现金流跟随利润增长 & CFO/NP、资本开支和融资成本 & CFO持续落后且融资成本上升 \\
\end{tabularx}
\sourcenote{analysis/implied\_growth\_sensitivity.md；analysis/risk\_framework.md}
\end{exhibitbox}
